\section{eFMI \& use-case workflow overview}\label{sec:efmi-and-use-case-workflow-overview}

This section provides a high-level overview of the eFMI Standard insofar as it is required to identify the main artifacts and generative model-based software engineering (MBSE) processes that it provides for our use-case; afterward, the objectives and workflow of the actual tooling we are using are presented.

\subsection{The eFMI Standard}\label{sec:efmi-standard}

The eFMI Standard\footref{fnt:efmi-website} \cite{eFMI_Standard} is an open standard of the Modelica Association\regtm{}\footnote{\tiny{\url{https://modelica.org/association/}}}. The official website\footref{fnt:efmi-website} provides a concise summary of the primary objectives of eFMI:
\blockquote{The \textit{eFMI\regtm{} Standard} is an open standard for the stepwise, model-transformation-based development of advanced control functions suited for safety-critical and real time targets. Its container architecture defines the common ground for collaboration among the stake-holders and toolchains along the various abstraction levels from high-level modeling and simulation -- e.g., a-causal, equation-based physics in Modelica -- down to actual embedded code. The model representations it defines are interlinked for traceability and semi-automatic code generation and do not only capture functional, but also non-functional quality requirements like coding standards, static memory allocation, worst time execution, and documentation, enabling credible development within a standardized workspace.\footnote{Originally written by Christoff Bürger.}}

\noindent{}The abstraction levels eFMI supports---each represented by a dedicated container type with well-defined content and with all containers of an embedded solution wrapped in a single embedded functional mock-up unit (eFMU)---are\footnote{Parts of the following Algorithm Code container section are from \cite{Kamp2025} and have been originally written by Christoff Bürger.} (note that, in contrast to FMUs, an eFMU is an open workspace container potentially still in development; eFMI can be used as a credible MBSE process to develop---i.e., implement---co-simulation FMUs for safety-critical real-time simulation):

\paragraph{Algorithm Code container:} Algorithmic high-level solution of a sampled system. Algorithms are implemented in GALEC, a new \underline{\textbf{g}}uarded \underline{\textbf{a}}lgorithmic \underline{\textbf{l}}anguage for \underline{\textbf{e}}mbedded \underline{\textbf{c}}ontrol defined in the eFMI Standard. Although imperative, GALEC provides many high-level concepts like multi-dimensional arithmetic operators, derived but well-defined dimensional sizes, and a built-in function standard library to abstract target-specific characteristics and provide common math operations, including interpolation and solving systems of linear equations. GALEC is computationally safe; a powerful static evaluation concept is used to guarantee all indexing is within bounds, all loops have an upper-bound, and all potential runtime errors like IEEE 754 not-a-number (NaN) values, singular linear equation systems or domain errors of built-in functions (like poles of trigonometric functions) are handled or explicitly exposed to the runtime environment (error signals of interface functions). GALEC is also well-defined regarding its imperative side-effects; for every two GALEC statements it is decidable whether they can be swapped, enabling, together with its well-defined dimensionality analysis, target-specific concurrent execution optimizations. Note that GALEC has no foreign function interface (FFI), only built-in functions; GALEC programs are therefore by nature self-contained. From a Modelica perspective, for example, a GALEC solution of a sampled system of equations (potentially with events or mixed system of equations) requires the solver to be inline-integrated as GALEC code, eliminating dependencies on third party numerical libraries.

These characteristics---self-contained, target independent, high abstraction level, and well-defined behavior under all circumstances---make GALEC an ideal intermediate language between simulation tools (that can target GALEC for code generation) and the embedded domain (where GALEC can serve as a source for code generation). Once a simulation tool finds a sampled solution for the simulation problem---which is not always possible for general equation systems---and expresses it as GALEC program, further tooling towards embedded C code "only" has to preserve these characteristics and can focus on code optimization (e.g., loop and multi-dimensional expression unrolling, or concurrent execution) and embedded target environment compliance (e.g., compliance with style guidelines, code analysers, interfaces, or other system integration requirements).

The sampled block-interface of the GALEC program of an Algorithm Code container is also exposed by the XML manifest file of the container, such that derived eFMU artifacts, like production code, can refer to it; this is important for traceability, but also to enable tooling to work with eFMUs without having to process GALEC (e.g., if just using production code without generating it). Naturally, each eFMU must have a single Algorithm Code container.

\paragraph{Production Code containers:} Target-specific production code (C program) implementing the GALEC program of the eFMU's Algorithm Code container. Typical challenges of production code generators are related to efficiently translating GALEC multi-dimensional arithmetic operations, providing implementations of the used built-in functions (or preferable reuse target ecosystem routines considering restricted availability of the C Standard Library), complying with embedded coding standards like MISRA C:2023 or the SEI CERT C Coding Standard\footnote{\label{fnt:sei-cert-c}\tiny{\url{https://cmu-sei.github.io/secure-coding-standards/sei-cert-c-coding-standard/}}}, following target ecosystem system integration APIs, generating ISO/IEC 9899 (C11, C17, C24, etc.) compliant code, and finding good trade-offs between portability and optimization.

An XML manifest file within each Production Code container maps the respective code's API to the well-defined (i.e., standardized in eFMI) block-interface of the eFMU's Algorithm Code container. This enables automatic system integration---as for example exploited by Dymola's production code export options presented in Section \ref{sec:production-code-export}---although eFMI does not standardize any production code interface (unlike FMI\footnote{\label{fnt:fmi}\tiny{\url{https://fmi-standard.org/}}} 3.0, which does with \texttt{fmi3Functions.h}). Naturally, an eFMU can contain many different production codes, e.g., for various floating-point precisions or code optimizations.

\paragraph{Binary Code containers:} Target specific binary code generated from one of the eFMU's Production Code containers. An XML manifest file lists the target platform, used compiler and linker options, and links back to the source Production Code container. Naturally, an eFMU can contain many different binaries, even for the same production code (e.g., based on different compiler and linker settings like debug and release versions, or for different targets).

\paragraph{Behavioral Model containers:} Test scenarios with reference trajectories stored in comma separated value (CSV) files and additional meta-information---like trajectory units, the index-wise splicing of multi-dimensional trajectories into separate scalar trajectories, absolute and relative tolerances, or scenario descriptions---stored in an XML manifest file. Two kinds of test scenarios are supported: variable time grid scenarios and equidistant time grid scenarios. The latter enforce a fixed sampling period within strict time tolerances between adjacent trajectory points, as, for example, typical for HiL tests. Variable time grid scenarios are more flexible, but require well-defined interpolation (and for that reason typically higher tolerances) when used for testing a sampled system like GALEC programs or derived production code; they typically stem from continuous simulations with flexible step size (e.g., DASSL).

Naturally, an eFMU can contain many Behavioral Model containers, e.g., created by different simulation or measurement tools. Automatic instantiation---although the APIs of production and binary codes are not standardized---for testing is possible thanks to all eFMU containers (including Behavioral Model ones) linking back to the well-defined block-interface of an eFMU's Algorithm Code container.

\subsection{Use-case tooling \& workflow}

In our use-case, we develop two models with the final objective of embedded deployment for CHiL testing: a power system stabilizer (i.e., an actual controller) and a hard real-time simulation of the controlled power system (i.e., the plant model). To that end, we used the following tools:

\squishlist
    \item \textbf{Dymola\regtm{}\footref{fnt:dymola} 2026x Refresh 1:} Physics-based design of Modelica models of the controller and plant; systematic MiL testing of design variants from continuous to sampled experiments with controller and plant models refined for eFMI (cf. Section \ref{sec:modeling-for-mil-testing}); controller and plant eFMU generation (cf. Section \ref{sec:eFMU-configurations}) and SiL tests of their production code against MiL experiments (cf. Section \ref{sec:modelica-based-sil-tests}).
    \item \textbf{Cppcheck\regtm{} Premium\footnote{\label{fnt:cppcheck-premium}\tiny{\url{https://www.cppcheck.com/}}} 26.3.0 and clang-tidy\footnote{\label{fnt:clang-tidy}\tiny{\url{https://clang.llvm.org/extra/clang-tidy/}}} (Clang Tools 19.1.5):} Strict code quality analyses---in particular for MISRA C:2023 compliance---of the eFMU production code of controller and plant (cf. Section \ref{sec:production-code-quality-analyses}).
    \item \textbf{STM32CubeMX\footref{fnt:stm32cubemx} 6.17.0:} Selection of STM32 microcontrollers for controller and plant; configuration of their peripherals---like pins, DACs, and timers---and generation of respective template code linking to the STM32CubeH7\footnote{\tiny{\url{https://www.st.com/en/embedded-software/stm32cubeh7.html}}} 1.13.0 software package with hardware abstraction layer (HAL), lower-layer (LL) drivers, and Arm\regtm{} Cortex\regtm{}-M Common Microcontroller Software Interface Standard (CMSIS) implementations\footnote{STM32CubeH7 also includes middleware libraries supporting, for example, a USB stack, TCP/IP networking stack, file system, or even real-time operation system (RTOS); we do not use any of the STM32CubeH7 middleware libraries.} (cf. Section \ref{sec:STM32CubeMX}).
    \item \textbf{STM32CubeIDE\footref{fnt:stm32cubeide} 2.1.0:} Import of CubeMX generated template codes of controller and plant as individual C development projects; incorporation of their eFMU Production Code containers as symbolically linked source and include directories; system integration of eFMU production codes into the \texttt{main()} execution routine of the microcontrollers; compilation and deployment on STM32 boards (cf. Section \ref{sec:STM32CubeIDE}).
    \item \textbf{WaveForms\footnote{\label{fnt:waveforms}\tiny{\url{https://digilent.com/shop/waveforms/}}} 3.25.1:} Interactive stimuli generation and plotting of CHiL test signals with a Digilent\regtm{} Analog Discovery 3 oscilloscope (cf. Section \ref{sec:rtchil}).
\squishend

\noindent{}Note, that our approach to deploy the plant model---which is required for CHiL testing the controller---directly on a microcontroller avoids the need for a real-time HiL platform, e.g., from dSPACE\regtm{}, Typhoon HIL\regtm{} or OPAL-RT Technologies\regtm{}; this reduces not only hardware costs but also the complexity of the required software and tooling stack and its licensing. All tools, except Dymola and Cppcheck Premium, are freely available. To generate eFMUs, the \textit{Dymola Source Code Generation} option, besides a standard Dymola license, is required. Although an open-source version of Cppcheck exists (and could have been used), its MISRA C:2023 and SEI CERT C Coding Standard\footref{fnt:sei-cert-c} coverage is significantly inferior to the commercial Premium version. And while production code analyses are optional for our use-case, they are of general interest for certification (Cppcheck Premium for example has TÜV SÜD\regtm{}\footnote{\tiny{\url{https://www.tuvsud.com/en}}} certification against EN 50128, EN 50657, EN 50716, IEC 61508, and ISO 26262).
